\documentclass[11pt]{article}
\usepackage[margin=1in]{geometry}
\usepackage{amsmath,amssymb,amsthm,hyperref}
\newtheorem{lemma}{Lemma}\newtheorem{proposition}[lemma]{Proposition}\newtheorem{corollary}[lemma]{Corollary}
\title{Integer-valued functional calculus for the normalized raising operator}
\author{Internal research note}\date{18 November 2026}
\begin{document}\maketitle
Fix $N\ge1$, $c=N+5$, and
$e_a=\lceil(s_7(N)-a-1)/6\rceil$.
\href{GAP_SIX_DIGIT_NORMALIZATION_AND_FINITE_GAUGE.tex}{The digit normalization theorem}, \texttt{thm:content}, gives integral columns
$\widetilde V_a(Q)=7^{-e_a}V_{N,a}(Q)$ for $Q\in\mathbb Z_7[X]$.
Let $\delta=vD_v$ and
\[
 D=\operatorname{diag}(c-a-6\delta),\quad L=D+J,
\]
where $J_{a,a-1}=a7^{e_{a-1}-e_a}$ for $1\le a<6$ and
$J_{0,5}=7^{1+e_5-e_0}v/5!$.
The exact raising identities in the proof of that theorem give
$\widetilde V(XQ)=L\widetilde V(Q)$.

\begin{lemma}[An integral formal gauge]\label{lem:formal}
One has $J^6=7vI$ and $[D,J]=-J$.
Both formal matrices
\[
 E_\pm=\sum_{j\ge0}\frac{(\pm J)^j}{j!}
          \in\operatorname{GL}_6(\mathbb Z_7[[v]])
\]
exist and are inverse, and $L=E_+DE_-$ on formal series.
No assertion that $E_\pm$ converge at $|v|_7=1$ is made.
\end{lemma}
\begin{proof}
The six link exponents sum to one, and the link constants
multiply to $(1\cdot2\cdot3\cdot4\cdot5)/5!=1$.
This gives $J^6=7vI$.
The internal diagonal difference is $-1$.
For the wrap link the Euler derivative contributes $-6$,
and the diagonal difference contributes $5$, again totaling $-1$.

Write $j=6\ell+r$, $0\le r<6$.
Each nonzero entry of $J^j$ has a seven-adic valuation
$\ell+\epsilon$ with $\epsilon\in\{0,1\}$.
Legendre's formula gives
\[
 v_7(J^j/j!)=\epsilon+\frac{s_7(j)-r}{6}\ge0.
\]
For $r>0$, $s_7(j)\ge r$; for $r=0,j>0$ its positive digit sum
is at least six. Each fixed $v$-coefficient receives only
finitely many $j$, so the two exponential series are integral
formal matrices. Their usual product and conjugation identities
are valid formally; alternatively multiply their finite
coefficients and use $[J,D]=J$.
The gauge is generally not restricted: choose any surviving
link of valuation zero and let $j=7^r$. Then
$J^j=(7v)^{(j-1)/6}J$ and
$v_7(j!)=(j-1)/6$, so that entry of $J^j/j!$
has a unit coefficient at arbitrarily high $v$-degree.
No cancellation with another $j$ occurs in that entry and degree.
The constant matrix is triangular unipotent at $v=0$,
so the formal inverses exist as stated.
\end{proof}

\begin{proposition}[All integer-valued polynomials preserve the lattice]
\label{prop:int}
If $Q\in\mathbb Q_7[X]$ takes $\mathbb Z_7$ into $\mathbb Z_7$,
then $Q(L)$ preserves $\mathbb Z_7[v]^6$ and
$\mathbb Z_7\langle v\rangle^6$.
In particular the canonical row-content bound remains valid
for every such $Q$:
\[
                  V_{N,a}(Q;v)\in7^{e_a}\mathbb Z_7[v].
\]
Moreover, on the Tate algebra the operator norm obeys
$\|Q(L)\|\le\sup_{x\in\mathbb Z_7}|Q(x)|_7$.
\end{proposition}
\begin{proof}
For every $m\ge0$, $\binom Dm$ acts coefficientwise on row $a$,
degree $k$, by the integer $\binom{c-a-6k}{m}$.
It therefore preserves $\mathbb Z_7[[v]]^6$.
The formal gauge shows that
$\binom Lm=E_+\binom Dm E_-$ does so too.
Acting on a polynomial it is a finite differential-polynomial
operator and has polynomial output. Hence that output is
integral; this argument never evaluates $E_\pm$ on the unit circle.

Every integer-valued polynomial has the finite Newton expansion
$Q(X)=\sum_m b_m\binom Xm$, with
$b_m=\Delta^mQ(0)\in\mathbb Z_7$.
Also $\max_m|b_m|_7=\sup_{\mathbb Z_7}|Q|_7$:
one direction follows from the difference formula and the other
from the integer-valued binomial expansion.
Thus $Q(L)$ has the asserted bound on polynomial vectors.
Their density extends the operator boundedly to the Tate algebra.
Apply it to the actual integral seed
$\widetilde V(1)$ and use the exact raising identity.
\end{proof}

\begin{corollary}[Extra divisibility for integer-valued test polynomials]
\label{cor:extra-content}
For $m\ge0$ and $Q\in\mathbb Q_7[X]$ taking $\mathbb Z_7$ into
$\mathbb Z_7$, one has
\[
 V_{N,a}((X^7-X)^mQ;v)\in 7^{e_a+m}\mathbb Z_7[v].
\]
\end{corollary}
\begin{proof}
The polynomial $(X^7-X)/7$ is integer-valued on $\mathbb Z_7$,
by Fermat's congruence. Apply Proposition~\ref{prop:int} to its
$m$th power times $Q$, then multiply by $7^m$.
\end{proof}

\begin{proposition}[Compatible residue projectors]\label{prop:projectors}
On $\mathbb Z_7\langle v\rangle^6$ the exact operator $L$ admits
a contractive functional calculus for continuous
$\mathbb Z_7$-valued functions on $\mathbb Z_7$.
For each $r\ge1$ the characteristic functions of the residue
classes modulo $7^r$ give bounded orthogonal projectors
summing to the identity. The projectors agree under refinement
of the residue classes.
\end{proposition}
\begin{proof}
For a residue class $b\bmod7^r$ with digits $b_i$, put
\[
 \rho_{b,r}(X)=
 \prod_{i=0}^{r-1}
 \left[1-\left(\binom X{7^i}-b_i\right)^6\right].
\]
Lucas' identity and continuity give $\rho_{b,r}(x)\equiv1$
on the class and zero off it. Thus $\rho_{b,r}(x)^{7^m}$ converges
uniformly to its characteristic function. On the class its
difference from one is divisible by $7^{m+1}$. Off the class
it is divisible by $7^{7^m}$. Consequently
$\rho_{b,r}(L)^{7^m}$ converges in operator norm by
Proposition~\ref{prop:int}.
Finite linear combinations of the characteristic functions
approximate every continuous $\mathbb Z_7$-valued function
uniformly, since the domain $\mathbb Z_7$ is compact.
The polynomial calculus therefore extends contractively to
these continuous functions. Multiplicativity passes to the
uniform limits, giving the claimed projector identities.
\end{proof}

The projectors act on the whole polynomial or Tate algebra
column. After evaluation at a unit parameter, components of the
actual seed can cancel. Controlling that cancellation is a
remaining step in the physical gap-six problem.
\end{document}
